% GENERATED FILE. Edit canonical lessons, then run npm run generate:books.
% canonical-source-hash: fnv1a64:f9fa83dbe29d686c
% canonical-lessons: HI-C309-rahat, HI-C309-sivir, HI-C309-dera, HI-C309-pirit, HI-C309-saran

\chapter{The Relief Camp}
\label{ch:hi-the-relief-camp}
\hlchaptermodality{309}

\begin{chapteropening}
\textbf{By the end of this chapter:} \emph{I can describe life in a relief camp.}

Complete the last lesson of chapter 309: I can describe life in a relief camp.
\end{chapteropening}

\section[\texorpdfstring{\hi{राहत} (rāhat)}{राहत (rāhat)}]{\hi{राहत} (rāhat) — relief}

\label{lesson:HI-C309-rahat}



\begin{quote}
Before the new one: say the Hindi for to set free, then the Hindi for to dart.
\end{quote}

\subsection*{You'll want to know: \hi{राहत}}
\textbf{\hi{राहत}} — \emph{rāhat} — \textquotedblleft{}relief\textquotedblright{}.

\begin{quote}
\hi{सरकार} \hi{ने} \hi{गाँव} \hi{में} \hi{राहत} \hi{भेजी।}
\end{quote}

\begin{quote}
\emph{sarkār ne gā̃v meṁ rāhat bhejī.}
\end{quote}

\begin{quote}
The government sent relief to the village.
\end{quote}

Feminine.

\begin{cousinweb}
\textbf{\hi{राहत}} (\emph{rāhat}) is usually traced, through Persian, to Arabic \emph{rāḥa}, \textquotedblleft{}rest, ease\textquotedblright{}. It is both the feeling of relief and the food and blankets sent after a disaster.
\end{cousinweb}

\begin{grammarlens}[title={What you've built}]
One more word for telling what happened.
\end{grammarlens}

\subsection*{Your turn}
\begin{itemize}
\raggedright
  \item \emph{Say it:} \emph{rāhat}
  \item \emph{Say it:} \emph{rāhat}, once more
  \item \emph{Say it:} \emph{sarkār ne gā̃v meṁ rāhat bhejī.} — The government sent relief to the village.
  \item \emph{Recall:} say the Hindi for to set free, then the Hindi for to dart, then say \emph{rāhat} again
\end{itemize}

\begin{tcolorbox}[breakable,colback=teal!4,colframe=teal!35!black,title={Before you move on}]
What does \hi{राहत} mean? (\textquotedblleft{}Relief\textquotedblright{}.) Say it once more.
\end{tcolorbox}
\section[\texorpdfstring{\hi{शिविर} (śivir)}{शिविर (śivir)}]{\hi{शिविर} (śivir) — a camp}

\label{lesson:HI-C309-sivir}



\begin{quote}
Before the new one: say the Hindi for to dart, then the Hindi for relief.
\end{quote}

\subsection*{You'll want to know: \hi{शिविर}}
\textbf{\hi{शिविर}} — \emph{śivir} — \textquotedblleft{}a camp\textquotedblright{}.

\begin{quote}
\hi{माफ़} \hi{कीजिए}, \hi{राहत} \hi{शिविर} \hi{कैसे} \hi{जाते} \hi{हैं}?
\end{quote}

\begin{quote}
\emph{māf kījiye, rāhat śivir kaise jāte haiṁ?}
\end{quote}

\begin{quote}
Excuse me, how do I get to the relief camp?
\end{quote}

Masculine.

\begin{cousinweb}
\textbf{\hi{शिविर}} (\emph{śivir}) is Sanskrit \emph{śibira} (also written \emph{śivira}), \textquotedblleft{}a royal camp, a tent\textquotedblright{} (Monier-Williams). Today it is a camp of any kind: \textbf{\hi{राहत} \hi{शिविर}}, a relief camp.
\end{cousinweb}

\begin{grammarlens}[title={What you've built}]
One more word for telling what happened.
\end{grammarlens}

\subsection*{Your turn}
\begin{itemize}
\raggedright
  \item \emph{Say it:} \emph{śivir}
  \item \emph{Say it:} \emph{śivir}, once more
  \item \emph{Say it:} \emph{māf kījiye, rāhat śivir kaise jāte haiṁ?} — Excuse me, how do I get to the relief camp?
  \item \emph{Recall:} say the Hindi for to dart, then the Hindi for relief, then say \emph{śivir} again
\end{itemize}

\begin{tcolorbox}[breakable,colback=teal!4,colframe=teal!35!black,title={Before you move on}]
What does \hi{शिविर} mean? (\textquotedblleft{}A camp\textquotedblright{}.) Say it once more.
\end{tcolorbox}
\section[\texorpdfstring{\hi{डेरा} (ḍerā)}{डेरा (ḍerā)}]{\hi{डेरा} (ḍerā) — an encampment; a temporary shelter}

\label{lesson:HI-C309-dera}



\begin{quote}
Before the new one: say the Hindi for relief, then the Hindi for a camp.
\end{quote}

\subsection*{You'll want to know: \hi{डेरा}}
\textbf{\hi{डेरा}} — \emph{ḍerā} — \textquotedblleft{}an encampment; a temporary shelter\textquotedblright{}.

\begin{quote}
\hi{बाढ़} \hi{के} \hi{बाद} \hi{लोग} \hi{नदी} \hi{के} \hi{किनारे} \hi{डेरे} \hi{में} \hi{रहे।}
\end{quote}

\begin{quote}
\emph{bāṛh ke bād log nadī ke kināre ḍere meṁ rahe.}
\end{quote}

\begin{quote}
After the flood people lived in a camp on the riverbank.
\end{quote}

Masculine.

\begin{cousinweb}
\textbf{\hi{डेरा}} (\emph{ḍerā}) is inherited from \emph{*ḍēra}, a form CDIAL reconstructs (CDIAL 5564), which gathers words for tents and resting places. \textbf{\hi{डेरे} \hi{में} \hi{रहना}} is to live in a camp or a shelter, away from home.
\end{cousinweb}

\begin{grammarlens}[title={What you've built}]
One more word for telling what happened.
\end{grammarlens}

\subsection*{Your turn}
\begin{itemize}
\raggedright
  \item \emph{Say it:} \emph{ḍerā}
  \item \emph{Say it:} \emph{ḍerā}, once more
  \item \emph{Say it:} \emph{bāṛh ke bād log nadī ke kināre ḍere meṁ rahe.} — After the flood people lived in a camp on the riverbank.
  \item \emph{Recall:} say the Hindi for relief, then the Hindi for a camp, then say \emph{ḍerā} again
\end{itemize}

\begin{tcolorbox}[breakable,colback=teal!4,colframe=teal!35!black,title={Before you move on}]
What does \hi{डेरा} mean? (\textquotedblleft{}An encampment; a temporary shelter\textquotedblright{}.) Say it once more.
\end{tcolorbox}
\section[\texorpdfstring{\hi{पीड़ित} (pīṛit)}{पीड़ित (pīṛit)}]{\hi{पीड़ित} (pīṛit) — a victim}

\label{lesson:HI-C309-pirit}



\begin{quote}
Before the new one: say the Hindi for a camp, then the Hindi for an encampment.
\end{quote}

\subsection*{You'll want to know: \hi{पीड़ित}}
\textbf{\hi{पीड़ित}} — \emph{pīṛit} — \textquotedblleft{}a victim\textquotedblright{}.

\begin{quote}
\hi{पीड़ितों} \hi{को} \hi{खाना} \hi{और} \hi{कंबल} \hi{मिले।}
\end{quote}

\begin{quote}
\emph{pīṛitoṁ ko khānā aur kambal mile.}
\end{quote}

\begin{quote}
The victims were given food and blankets.
\end{quote}

Masculine or feminine, by the person meant.

\begin{cousinweb}
\textbf{\hi{पीड़ित}} (\emph{pīṛit}) is Sanskrit \emph{pīḍita}, \textquotedblleft{}pressed, afflicted\textquotedblright{} (Monier-Williams). As a noun it is the person who suffered: \textbf{\hi{बाढ़} \hi{पीड़ित}}, flood victims.
\end{cousinweb}

\begin{grammarlens}[title={What you've built}]
One more word for telling what happened.
\end{grammarlens}

\subsection*{Your turn}
\begin{itemize}
\raggedright
  \item \emph{Say it:} \emph{pīṛit}
  \item \emph{Say it:} \emph{pīṛit}, once more
  \item \emph{Say it:} \emph{pīṛitoṁ ko khānā aur kambal mile.} — The victims were given food and blankets.
  \item \emph{Recall:} say the Hindi for a camp, then the Hindi for an encampment, then say \emph{pīṛit} again
\end{itemize}

\begin{tcolorbox}[breakable,colback=teal!4,colframe=teal!35!black,title={Before you move on}]
What does \hi{पीड़ित} mean? (\textquotedblleft{}A victim\textquotedblright{}.) Say it once more.
\end{tcolorbox}
\section[\texorpdfstring{\hi{शरण} (śaraṇ)}{शरण (śaraṇ)}]{\hi{शरण} (śaraṇ) — refuge; shelter}

\label{lesson:HI-C309-saran}



\begin{quote}
Before the new one: say the Hindi for an encampment, then the Hindi for a victim.
\end{quote}

\subsection*{You'll want to know: \hi{शरण}}
\textbf{\hi{शरण}} — \emph{śaraṇ} — \textquotedblleft{}refuge; shelter\textquotedblright{}.

\begin{quote}
\hi{पीड़ितों} \hi{ने} \hi{पास} \hi{के} \hi{मंदिर} \hi{में} \hi{शरण} \hi{ली।}
\end{quote}

\begin{quote}
\emph{pīṛitoṁ ne pās ke mandir meṁ śaraṇ lī.}
\end{quote}

\begin{quote}
The victims took refuge in a nearby temple.
\end{quote}

Feminine.

\begin{cousinweb}
\textbf{\hi{शरण}} (\emph{śaraṇ}) is Sanskrit \emph{śaraṇa}, \textquotedblleft{}a shelter, a protection\textquotedblright{} (Monier-Williams), feminine in Hindi. \textbf{\hi{शरण} \hi{लेना}} is to take refuge.
\end{cousinweb}

\begin{grammarlens}[title={What you've built}]
One more word for telling what happened.
\end{grammarlens}

\subsection*{Your turn}
\begin{itemize}
\raggedright
  \item \emph{Say it:} \emph{śaraṇ}
  \item \emph{Say it:} \emph{śaraṇ}, once more
  \item \emph{Say it:} \emph{pīṛitoṁ ne pās ke mandir meṁ śaraṇ lī.} — The victims took refuge in a nearby temple.
  \item \emph{Recall:} say the Hindi for an encampment, then the Hindi for a victim, then say \emph{śaraṇ} again
\end{itemize}

\begin{tcolorbox}[breakable,colback=teal!4,colframe=teal!35!black,title={Before you move on}]
What does \hi{शरण} mean? (\textquotedblleft{}Refuge; shelter\textquotedblright{}.) Say it once more.
\end{tcolorbox}
